% =====================================================================
% Section: Experiments -- Datasets and implementation details (outline.md §3.1;
%   defines sec:exp_setup). Opens the Experiments section.
% Draft v1 (2026-06-17). All settings are transcribed verbatim from the training
%   configs (CATANet/experiments/*/*.yml) and the efficiency/test logs recorded in
%   paper/plan/data-collection-checklist.md (A2/A3) and experiment-protocol.md.
% This draft makes NO result assertions; it only states the protocol.
% HARDWARE (2026-06-28): per the user's instruction we adopt the same hardware as
%   the CATANet paper. CATANet only reports running each model on a single
%   NVIDIA RTX 4090 GPU (CATANet §4.5 efficiency measurement); training GPU count
%   and total GPU-hours are not disclosed in CATANet; they are not reported.
%   Both training and efficiency measurement therefore state a single RTX 4090.
% =====================================================================
\section{Experiments}
\label{sec:experiments}

\subsection{Datasets and implementation details}
\label{sec:exp_setup}

\paragraph{Datasets and metrics.}
Following the standard lightweight-SR protocol, we train on the $800$ training
images of DIV2K~\cite{agustsson2017div2k} and evaluate on five benchmarks:
Set5~\cite{bevilacqua2012set5}, Set14~\cite{zeyde2010set14},
BSD100~\cite{martin2001bsd100}, Urban100~\cite{huang2015urban100}, and
Manga109~\cite{matsui2017manga109}. Low-resolution inputs are produced by
bicubic downscaling at $\times2$, $\times3$, and $\times4$. We report PSNR and
SSIM~\cite{wang2004ssim} on the Y channel of the YCbCr space, discarding a border
of $s$ pixels for scale $s$, which matches the convention used by the cited
baselines.

\paragraph{Network configuration.}
DPRNet keeps the CATANet backbone unchanged and replaces only the routing inside
each TAB with the DPR module (Section~\ref{sec:method}). The deep feature extractor
stacks $L{=}8$ blocks with channel width $C{=}40$. The number of prototype slots
per block is $M=[16,32,64,128,16,32,64,128]$, and the intra-group attention uses a
group size of $128$. The learnable routing temperature is initialized at $\tau{=}6$
(i.e.\ a log-scale parameter $\theta{=}\log 6$) and clamped at $\tau_{\max}{=}10$
(Eq.~\eqref{eq:scores}); the per-block usage-balancing weights are
$\lambda=[5{\times}10^{-4}]_{\times6}$ for the first six blocks and
$7{\times}10^{-4}$ for the last two (Eq.~\eqref{eq:balance}). These yield
$0.60$--$0.67$\,M parameters depending on scale (Table~\ref{tab:efficiency}).

\paragraph{Training.}
We train the $\times2$ model for $800$k iterations and the $\times3$ and $\times4$
models for $250$k iterations each. Training patches are $128{\times}128$, $192{\times}192$,
and $256{\times}256$ HR crops for $\times2/\times3/\times4$ respectively, with a
batch size of $16$ and random horizontal flips and $90^{\circ}$ rotations for
augmentation. We optimize the $\ell_1$ reconstruction loss together with the
usage-balancing regularizer (Eq.~\eqref{eq:balance}) using Adam
($\beta_1{=}0.9$, $\beta_2{=}0.99$, no weight decay) at an initial learning rate of
$2{\times}10^{-4}$, halved by a MultiStep schedule (milestones at
$\{300,500,650,700,750\}$k for $\times2$ and $\{125,200,225,237.5\}$k for
$\times3/\times4$). The learnable routing temperature is trained with a
$0.1{\times}$ learning-rate multiplier for stability. Validation is run every
$2{.}5$k iterations. Experiments use PyTorch with the BasicSR
framework on a single NVIDIA RTX~4090 GPU.

\paragraph{Efficiency measurement.}
Parameters, Multi-Adds (MACs reported by thop), and inference latency are measured
on a single NVIDIA RTX~4090 GPU for a fixed LR input of $256{\times}256$ ($\times4$
output $1024{\times}1024$), averaged over $100$ runs after $20$ warm-up runs. This input
convention matches the one under which the CATANet paper reports Multi-Adds for
itself and the lightweight Transformer baselines, making the figures directly
comparable. Efficient-SR challenge reports likewise treat runtime, parameters, and
FLOPs as joint deployment criteria rather than auxiliary statistics~\cite{ren2024ntire}.
For comparison with FLOPs-based reports, Multi-Adds can be multiplied by two.
